\documentclass{article}
\usepackage{graphicx} % Required for inserting images
\usepackage{amsmath}
% Настройка шрифтов и поддержки русского языка
\usepackage[T2A]{fontenc}
\usepackage[utf8]{inputenc}
\usepackage[russian]{babel}

% Важно: объявляем оператор \rank, так как стандартной команды для него в LaTeX нет
\usepackage{amsmath}
\DeclareMathOperator{\rank}{rank} 

\title{Home work 1}
\author{Dmitri Churikov}
\date{September 2026}

\begin{document}
\maketitle

\section{Задание}

Определи ранг произведения DA и обоснуй вывод.
\\

Решение:
Согласно Теореме о ранге (документ Неделя 1 Методы оптимизации Лекция.pdf) 
\\
\rank(AB) \le \min\{\rank(A), \rank(B)\}

Если матрица  D невырождена ((\(\det A \neq 0\))), ранг произведения равен r

Если вырождена — ранг произведения уменьшается вплоть до нуля.


\section{Задание}
Дано:  \\
$Ax = \lambda x$\\
1) Найти: A + cI \\

Решение:

Умножим новую матрицу \(\mathbf{A} + c\mathbf{I}\) на тот же собственный вектор:

Получим:
 (A+cI)x=Ax+cIx
 
 Так как \(\mathbf{I}\mathbf{x} = \mathbf{x}\) 

 \[\mathbf{Ax}+c\mathbf{Ix}=\lambda \mathbf{x}+c\mathbf{x}\]

Вынесем вектор x за скобки:\[\lambda \mathbf{x}+c\mathbf{x}=(\lambda +c)\mathbf{x}\] 

Получим:
A + cI = \lambda + c \\


2) Найти A^k

Решение:
Умножим обе части равенства слева на матрицу A:

\[\mathbf{A}(\mathbf{Ax})=\mathbf{A}(\lambda \mathbf{x})\]

\[\mathbf{A}^{2}\mathbf{x}=\lambda (\mathbf{Ax})\]

\[\mathbf{A}^{2}\mathbf{x}=\lambda (\lambda \mathbf{x})=\lambda ^{2}\mathbf{x}\]

Повторяя эту операцию k раз, мы получаем:\[\mathbf{A}^{k}\mathbf{x}=\lambda ^{k}\mathbf{x}\]

Ответ: A^k = \lambda ^{k}


3)  \mathbf{A}^{-1} - ?

Решение:
Умножим обе части равенства слева на матрицу \mathbf{A}^{-1}:

\[\mathbf{A}^{-1}(\mathbf{Ax})=\mathbf{A}^{-1}(\lambda \mathbf{x})\]

В левой части получаем единичную матрицу: \(\mathbf{A}^{-1}\mathbf{A}\mathbf{x} = \mathbf{I}\mathbf{x} = \mathbf{x}\)

В правой части выносим скаляр \lambda \вперед: \(\lambda (\mathbf{A}^{-1}\mathbf{x})\)

\[\mathbf{x}=\lambda \mathbf{A}^{-1}\mathbf{x}\]

Разделим обе части на \lambda : 

\[\mathbf{A}^{-1}\mathbf{x}=\frac{1}{\lambda }\mathbf{x}\]

Ответ:
\mathbf{A}^{-1}  = \frac{1}{\lambda }

\section{Задание}
Задание 3.1: Вычисли сингулярное разложение SVD следующих матриц:

Матрица \(A_1 = \begin{pmatrix} 2 \\ 2 \\ 8 \end{pmatrix}\)

Матрица \Sigma (размер \(3 \times 1\)):

Найдём норму вектора: \(\sqrt{2^2 + 2^2 + 8^2} = \sqrt{4 + 4 + 64} = \sqrt{72} = 6\sqrt{2}\).
\[\Sigma =\left(\begin{matrix}6\sqrt{2}\\ 0\\ 0\end{matrix}\right)

Матрица V^{T} (размер \(1 \times 1\)):Поскольку это скалярное пространство для правого вектора:\[V^{T}=\left(\begin{matrix}1\end{matrix}\right)\]

Матрица U(размер \(3 \times 3\)):Первый столбец получается нормированием исходной матрицы A_{1}: \(\frac{1}{6\sqrt{2}} \begin{pmatrix} 2 \\ 2 \\ 8 \end{pmatrix} = \begin{pmatrix} \frac{\sqrt{2}}{6} \\ \frac{\sqrt{2}}{6} \\ \frac{2\sqrt{2}}{3} \end{pmatrix}\). Остальные два столбца дополняют матрицу до ортонормированного базиса.\[U=\left(\begin{matrix}\frac{\sqrt{2}}{6}&-\frac{\sqrt{2}}{2}&-\frac{2\sqrt{2}}{3}\\ \frac{\sqrt{2}}{6}&\frac{\sqrt{2}}{2}&-\frac{2\sqrt{2}}{3}\\ \frac{2\sqrt{2}}{3}&0&\frac{\sqrt{2}}{3}\end{matrix}\right)\approx \left(\begin{matrix}0.236&-0.707&-0.667\\ 0.236&0.707&-0.667\\ 0.943&0&0.333\end{matrix}\right)\]

Задание 3.2 (дата рождения 01/10):

Матрица \(A_2 = \begin{pmatrix} 0 & 11 \\ 11 & 0 \\ 0 & 0 \end{pmatrix}\)

 \(A_2^T A_2 = \begin{pmatrix} 121 & 0 \\ 0 & 121 \end{pmatrix}\). \[\Sigma =\left(\begin{matrix}11&0\\ 0&11\\ 0&0\end{matrix}\right)\]

Матрица V^{T}(размер \(2 \times 2\)):

В качестве базиса можно выбрать стандартную единичную матрицу:\[V^{T}=\left(\begin{matrix}1&0\\ 0&1\end{matrix}\right)\]

Матрица U (размер \(3 \times 3\)):\\

Первые два столбца формируются как \(u_i = \frac{1}{\sigma_i} A_2 v_i\), а третий столбец ортогонален им:\[U=\left(\begin{matrix}0&1&0\\ 1&0&0\\ 0&0&1\end{matrix}\right)\]
\end{document}
